%% ================================================================
%%  BODY SKELETON -- paste between \begin{document} and \end{document}
%%  in the Overleaf template, after the preamble block.
%%
%%  PAGE BUDGET (8 pp excluding references). Track it as we fill:
%%    1 Introduction ............ 1.00
%%    2 Related Work ............ 0.50
%%    3 Reconstruction .......... 1.25
%%    4 The MDP ................. 1.50
%%    5 Experiments ............. 2.50
%%    6 Where It Breaks ......... 1.00
%%    7 Limitations & Outlook ... 0.25
%%                        total = 8.00
%%
%%  Sections can be filled in ANY order. Every one is self-contained
%%  and carries a \label, so \cref works before its neighbours exist.
%% ================================================================

\maketitle

\begin{abstract}
\TODO{abstract --- write LAST, once section 5 is settled. Roughly 150 words:
the question (geometry vs. weights), the MDP in one sentence, the headline
number, the wall at long horizons and its named cause.}
\end{abstract}

% ================================================================
\section{Introduction}
\label{sec:intro}
% Target: 1.00 page. Write LAST but for the abstract -- it must promise
% exactly what section 5 delivers, no more.
% Beats: (i) world models are rolled out globally, physical AI needs
% query-driven local prediction; (ii) a scheme = geometry + weights, and
% only the geometry was ever a free choice; (iii) the question; (iv) why
% the answer is interesting either way; (v) contributions.
% ================================================================

\TODO{section 1}

\paragraph{Contributions.}
\TODO{3--4 bullets. Keep each one to a claim we can point at a figure for.}

% ================================================================
\section{Related Work}
\label{sec:related}
% Target: 0.50 page. Three \paragraph blocks, no subsections.
% MUST cite the RL-for-AMR line explicitly -- a reviewer who knows it and
% does not see it cited will assume we do not. State the three
% differentiators plainly: single query vs. global mesh; space--time vs.
% spatial mesh per step; consistency structural vs. learned.
% ================================================================

\paragraph{Neural surrogates for PDEs.}
\TODO{operator learning / physics-informed nets --- one or two sentences
positioning them as global-rollout models.}

\paragraph{Reinforcement learning for mesh adaptation.}
\TODO{the closest prior work. \CITE{yang2021 rl-amr}, \CITE{foucart2023 deep-rl-amr},
\CITE{freymuth2023 swarm-amr}. Then the three differentiators.}

\paragraph{Meshfree and generalised finite differences.}
\TODO{GFDM lineage --- establishes that the reconstruction operator is
standard, so the novelty is unambiguously the placement policy.}

% ================================================================
\section{Reconstruction: Consistency by Construction}
\label{sec:recon}
% Target: 1.25 pages. This is SETUP, not contribution -- it earns its space
% only by establishing that every scheme the agent can express is
% consistent. If the paper runs long, cut here first, never section 5.
% Equations lift verbatim from MAIN.tex ch. 2 (notation macros match).
% ================================================================

\subsection{Local linear prediction and the moment conditions}
\label{sec:recon_taylor}
\TODO{predictor $\uhat = \bw^\top \bu$; moment conditions; underdetermined
$n > m$; minimum-norm solution $\bw = \bA^+ \bb$.}

\subsection{Substituting the PDE into the expansion}
\label{sec:recon_pde}
\TODO{substitution vs. augmentation; the three-row constraint matrix;
order $p$ needs $p+1$ rows, not $(p+1)(p+2)/2$.}

\subsection{Error bound and the role of $\normone{\bw}$}
\label{sec:recon_bound}
\TODO{$|\ux - \uhat| \le \normone{\bw}\, C h^2$; positivity $\Rightarrow$
convex combination $\Rightarrow \normone{\bw} = 1$ $\Rightarrow$ linear
rather than geometric error accumulation. This is the hook that
section~\ref{sec:breaks} pays off --- set it up carefully here.}

% ================================================================
\section{Placement as a Markov Decision Process}
\label{sec:mdp}
% Target: 1.50 pages. The section an ML reviewer reads hardest.
% Lead with the framing: the actions are COMPUTATIONS, not physical motions.
% ================================================================

\subsection{State}
\label{sec:mdp_state}
\TODO{$\R^{2N+2}$: walker offsets relative to $\zstar$ in grid units, plus
$t^*/\Delta t$ and $x^*/\Delta x$. Say plainly that the visited set is NOT
observed --- it is a POMDP and we should own that, not hide it.}

\subsection{Actions and admissibility masking}
\label{sec:mdp_action}
\TODO{$\Aset$, the flattened $N \times |\Aset|$ discrete space, and why
masking was necessary: a rejected move leaves the observation unchanged,
so a deterministic policy re-selects it forever. Numerical admissibility
(two-sided coverage, conditioning) supplies the mask.}

\subsection{Termination}
\label{sec:mdp_term}
\TODO{the bubble; and the key design point --- the stopping rule asks the
same question as the scoring rule, i.e. "does an admissible stencil exist
at $\zstar$", not "are $n$ points nearby".}

\subsection{Reward}
\label{sec:mdp_reward}
\TODO{$-w_{\mathrm{err}} \enorm - w_{\mathrm{time}} (K/\Kmax)$, the log
normalisation of $\enorm$, and potential-based shaping. Worth stating that
the shaping is policy-invariant (telescopes to a constant), so it is a
variance knob and not a change of objective.}

\subsection{Policy architecture and training}
\label{sec:mdp_policy}
\TODO{flattened joint softmax and why the factored/MultiDiscrete form
failed; MaskablePPO, $\gamma = 1$, hyperparameters.}

% \begin{algorithm}[t]
% \caption{One episode of query-driven placement}
% \label{alg:episode}
% \begin{algorithmic}[1]
% \TODO{pseudocode}
% \end{algorithmic}
% \end{algorithm}

% ================================================================
\section{Experiments}
\label{sec:exp}
% Target: 2.50 pages. The load-bearing section. Nothing here may live in
% an appendix -- the workshop has no supplementary upload.
% ================================================================

\subsection{Setup}
\label{sec:exp_setup}
\TODO{viscous Burgers via Hopf--Cole; grid; what the agent is given
(initial condition and boundaries) and what it must compute (the interior).
Be exact about this --- it is the claim most likely to be probed.}

\subsection{Does the policy learn?}
\label{sec:exp_learn}
\TODO{learning curves: reached rate, terminal error, steps to termination.
The masked-random baseline goes here as the floor.}

% \begin{figure}[t]
%   \centering
%   \TODO{learning curves figure}
%   \caption{\TODO{caption}}
%   \label{fig:learning}
% \end{figure}

\subsection{Cost of a single query}
\label{sec:exp_workprecision}
\TODO{work--precision: terminal error vs. number of evaluated points, for
the learned policy, masked-random, and a uniform grid march, at matched
error. THE headline claim.}

% \begin{figure}[t]
%   \centering
%   \TODO{work-precision figure}
%   \caption{\TODO{caption}}
%   \label{fig:workprecision}
% \end{figure}

\subsection{What geometry does the policy discover?}
\label{sec:exp_geometry}
\TODO{the visited cloud; envelope fit $w \propto \tau^p$ to separate
diffusive ($p \approx 1/2$) from characteristic ($p \approx 1$) from
shortest-path artifact. Note: at short horizons $\enorm \equiv 0$, so the
tent shape there proves nothing --- use the horizon where accuracy is
actually active.}

% \begin{figure}[t]
%   \centering
%   \TODO{trajectory / envelope figure}
%   \caption{\TODO{caption}}
%   \label{fig:geometry}
% \end{figure}

\subsection{Ablations}
\label{sec:exp_ablation}
\TODO{pick from the experimental record: masking on/off, state version,
reward era, bubble radius. One compact table beats four small plots.}

\begin{table}[t]
  \centering
  \caption{\TODO{ablation caption}}
  \label{tab:ablation}
  \TODO{ablation table}
\end{table}

% ================================================================
\section{Where the Method Breaks}
\label{sec:breaks}
% Target: 1.00 page. The payoff of the framing -- do not bury this and do
% not apologise for it. We asked a question, got a partial yes, and can
% name the obstruction exactly.
% ================================================================

\subsection{The long-horizon wall}
\label{sec:breaks_wall}
\TODO{where and how it fails as $t^*$ grows.}

\subsection{The obstruction is stability, not learning}
\label{sec:breaks_why}
\TODO{the amplification argument: $\normone{\bw} > 1$ gives geometric
rather than linear error growth; a coordinated front gives
$\normone{\bw} = 1$, an isolated zigzagging walker does not. Emphasise
that conditioning does not detect this --- both geometries pass
$\mathrm{cond}(\bA) < 10^4$ comfortably.}

% ================================================================
\section{Limitations and Outlook}
\label{sec:conclusion}
% Target: 0.25 page. Be direct. Single initial condition, reward needs the
% exact solution, POMDP state. Then the one-sentence outlook.
% ================================================================

\TODO{section 7}

\begin{ack}
\TODO{funding / compute / acknowledgements. Not anonymised --- this
workshop is single-blind.}
\end{ack}

\bibliographystyle{plainnat}
\bibliography{refs}
% Create refs.bib in the Overleaf project. Nothing after this line counts
% toward the 8 pages.

% ================================================================
%  APPENDIX -- same PDF, after the references. There is NO supplementary
%  upload on OpenReview, so anything here must be strictly optional:
%  extra trajectory figures, the full experimental record table.
%  Never put a load-bearing result here.
% ================================================================
% \appendix
% \section{Full experimental record}
% \label{app:record}
